\documentclass[10pt,a4paper]{amsart}
\usepackage{amsmath,amssymb,amsfonts}
\usepackage{enumitem}
\usepackage[left=1.5cm, right=1.5cm]{geometry}

\newcommand{\R}{\mathbb{R}}
\renewcommand{\d}{\mathrm{d}}

\title{Introdução às Equações Diferenciais - Lista 4}

\begin{document}
\maketitle





\textbf{Problema 1.-} Seja \(f(x)\) uma função real tal que
\[
\lim_{x\to +\infty} \frac{f(x)}{e^x}=1
\]
e
\[
|f''(x)|<c|f'(x)|
\]
para todo \(x\) suficientemente grande. Prove que
\[
\lim_{x\to +\infty} \frac{f'(x)}{e^x}=1.
\]

\vspace{5mm}

\textbf{Problema 2.-} Sejam \(a(x)\) e \(r(x)\) funções positivas e contínuas definidas no intervalo
\([0,\infty)\), e suponha que
\[
\liminf_{x\to\infty} \bigl(x-r(x)\bigr)>0.
\]
Suponha que \(y(x)\) seja uma função contínua em toda a reta real, que seja
diferenciável em \([0,\infty)\), e que satisfaça
\[
y'(x)=a(x)y(x-r(x))
\]
em \([0,\infty)\). Prove que o limite
\[
\lim_{x\to\infty} y(x)\exp\left\{-\int_0^x a(u)\,du\right\}
\]
existe e é finito.

\vspace{5mm}

\textbf{Problema 3.-} Sejam \(l_0,c,\alpha,g\) constantes positivas, e seja \(x(t)\) a solução da
equação diferencial
\[
\bigl((l_0+ct^\alpha)^2x'\bigr)'
+g[l_0+ct^\alpha]\sin x=0,
\qquad t\geq 0,\qquad -\frac{\pi}{2}<x<\frac{\pi}{2},
\]
satisfazendo as condições iniciais
\[
x(t_0)=x_0,\qquad x'(t_0)=0.
\]
Esta é a equação do pêndulo matemático cujo comprimento varia de acordo
com a lei \(l=l_0+ct^\alpha\). Prove que \(x(t)\) está definida no intervalo
\([t_0,\infty)\); além disso, se \(\alpha>2\), então, para todo \(x_0\neq 0\),
existe \(t_0\) tal que
\[
\liminf_{t\to\infty}|x(t)|>0.
\]

\vspace{5mm}

\textbf{Problema 4.-} Prove que, se \(a_i\) \((i=1,2,3,4)\) são constantes positivas,
\(a_2-a_4>2\), e \(a_1a_3-a_2>2\), então a solução
\((x(t),y(t))\) do sistema de equações diferenciais
\[
\dot{x}=a_1-a_2x+a_3xy,
\]
\[
\dot{y}=a_4x-y-a_3xy
\qquad (x,y\in\mathbb{R})
\]
com as condições iniciais
\[
x(0)=0,\qquad y(0)\geq a_1
\]
é tal que a função \(x(t)\) possui exatamente um máximo local estrito
no intervalo \([0,\infty)\).

\vspace{5mm}

\textbf{Problema 5.-} Seja \(x:[0,\infty)\to\mathbb{R}\) uma função diferenciável que satisfaz a identidade
\[
x'(t)
=
-2x(t)\sin^2 t
+
(2-|\cos t|+\cos t)
\int_{t-1}^{t} x(s)\sin^2 s\,ds
\]
em \([1,\infty)\). Prove que \(x\) é limitada em \([0,\infty)\) e que
\[
\lim_{t\to\infty}x(t)=0.
\]
A conclusão continua verdadeira para funções que satisfazem a identidade
\[
x'(t)
=
-2x(t)
+
(2-|\cos t|+\cos t)
\int_{t-1}^{t}x(s)\,ds?
\]

\vspace{5mm}

\textbf{Problem 6.-} Sejam \(A,B\) matrizes reais ou complexas. Prove que
\[
e^{t(A+B)}=e^{tA}\cdot e^{tB}
\]
para todo \(t\in\mathbb{R}\) se e somente se
\[
AB=BA.
\]

\vspace{5mm}

\textbf{Problema 7.-}
Seja \(f:[0,\infty)\to\mathbb{R}\) uma função duas vezes diferenciável tal que
\[
f(0)=0,\qquad f'(0)=1,
\]
e
\[
f''(x)=\frac{1}{1+f(x)}
\]
para todo \(x\in[0,1]\). Prove que
\[
f(1)<\frac{3}{2}.
\]

\vspace{5mm}

\textbf{Problema 8.-}
Seja \(a>0\). Encontre todas as funções continuamente diferenciáveis
\[
\varphi:[a,\infty)\to\mathbb{R}
\]
tais que, para todo \(x>a\),
\[
\varphi(x)^2
=
\int_a^x \left(|\varphi(y)|^2+|\varphi'(y)|^2\right)\,dy
-(x-a)^3.
\]

\vspace{5mm}

\textbf{Problema 9.-}
Seja
\[
f_0:[0,1]\to\mathbb{R}
\]
uma função contínua. Para \(n\geq 1\), defina
\[
f_n(x)=\int_0^x f_{n-1}(t)\,dt.
\]
Prove que a série
\[
\sum_{n=1}^{\infty} f_n(x)
\]
converge para todo \(x\in[0,1]\), e encontre uma fórmula explícita para sua soma.

\vspace{5mm}

\textbf{Problem 10.-}
\begin{enumerate}
    \item Encontre a transformada de Laplace das seguintes funções.

    \begin{enumerate}
        \item 
        \[
        f(t)=\sin(2t)\cos(2t).
        \]

        \item 
        \[
        f(t)=\cos^2(3t).
        \]

        \item 
        \[
        f(t)=t e^{2t}\sin(3t).
        \]

        \item 
        \[
        f(t)=(t+3)u_7(t).
        \]

        \item 
        \[
        f(t)=t^2u_3(t).
        \]

        \item 
        \[
        f(t)=
        \begin{cases}
        1, & 0\leq t<2,\\
        t^2-4t+4, & t\geq 2.
        \end{cases}
        \]

        \item 
        \[
        f(t)=
        \begin{cases}
        t, & 0\leq t<3,\\
        5, & t\geq 3.
        \end{cases}
        \]

        \item 
        \[
        f(t)=
        \begin{cases}
        0, & t<\pi,\\
        t-\pi, & \pi\leq t<2\pi,\\
        0, & t\geq 2\pi.
        \end{cases}
        \]

        \item 
        \[
        f(t)=
        \begin{cases}
        \cos(\pi t), & t<4,\\
        0, & t\geq 4.
        \end{cases}
        \]

        \item 
        \[
        f(t)=
        \begin{cases}
        t, & 0\leq t<1,\\
        e^t, & t\geq 1.
        \end{cases}
        \]
    \end{enumerate}

    \item Encontre a transformada inversa de Laplace.

    \begin{enumerate}
        \item 
        \[
        F(s)=\frac{1}{(s+1)(s^2-1)}.
        \]

        \item 
        \[
        F(s)=\frac{2s+3}{s^2+4s+13}.
        \]

        \item 
        \[
        F(s)=\frac{e^{-3s}}{s-2}.
        \]

        \item 
        \[
        F(s)=\frac{1+e^{-2s}}{s^2+6}.
        \]
    \end{enumerate}
\end{enumerate}








\end{document}